\documentclass{beamer}
\usepackage{amssymb}


\begin{document}

  \begin{frame}
    \textit{Teorema}: $A \leq B$ si y sólo si $-B \leq -A$.
  \end{frame}

  \begin{frame}
    \textit{Teorema}: $A \leq B$ si y sólo si $-B \leq -A$.
    \newline\newline

    \textit{Teorema}: $A \in \mathbb{R}^+$ si y sólo si $-A \in \mathbb{R}^-$.
  \end{frame}

  \begin{frame}
    \textit{Definción}: Dados $A, B \subseteq \mathbb{Q}$, denotamos el 
    siguiente conjunto

    \[A \cdot B = \{x \cdot y \in \mathbb{Q}: x \in A \land y \in B\}\]
  \end{frame}

  \begin{frame}
    \textit{Teorema}: Si $A, B \in \mathbb{R}^+$ entonces $A \cdot B \in \mathbb{R}^+$.
  \end{frame}

  \begin{frame}
    \textit{Teorema}: Dados $A, B, C \in \mathbb{R}^+$:

    \begin{itemize}
      \item $AB = BA$.
      \item $A(BC) = (AB)C$.
      \item $A(B + C) = AB + AC$.
    \end{itemize}
  \end{frame}

  \begin{frame}
    \textit{Teorema}: Dados $A, B, C \in \mathbb{R}^+$:

    \begin{itemize}
      \item $AB = BA$.
      \item $A(BC) = (AB)C$.
      \item $A(B + C) = AB + AC$.
    \end{itemize}

    \textit{Teorema}: Para todo real $A$ positivo, $A \cdot 1 = A$.
  \end{frame}

  \begin{frame}
    \textit{Teorema}: Para todo par de racionales positivos $p$ y $q$, se 
    satisface que:

    \[S^{>_\mathbb{Q}}(p) S^{>_\mathbb{Q}}(q) = S^{>_\mathbb{Q}}(pq)\]
  \end{frame}

  \begin{frame}
    \textit{Definción}: Dado $A \in \mathbb{R}^+$, se define 
    
    \begin{equation*}
      A^{-1} = \left\{
      \begin{aligned}
        &S^{>_\mathbb{Q}}(\frac{1}{r}) \text{ si } A = S^{>_\mathbb{Q}}(r), \\
        &{x^{-1}: x \in A^c \cap \mathbb{Q}^+} \text{ si } A \not\in \mathbb{Q}
      \end{aligned}
      \right.
    \end{equation*}
  \end{frame}

  \begin{frame}
    
    \textit{Teorema}: Si $A \in \mathbb{R}^+$, entonces $A A^{-1} = 1$.
  \end{frame}

  \begin{frame}
    \textit{Definición}: Definimos la muliplicación en $\mathbb{R}$ de la siguiente forma

    \begin{equation*}
      AB = \left\{
      \begin{aligned}
        &AB \text{ si } A, B > 0
      \end{aligned}
      \right.
    \end{equation*}
  \end{frame}


  \begin{frame}
    \textit{Definición}: Definimos la muliplicación en $\mathbb{R}$ de la siguiente forma

    \begin{equation*}
      AB = \left\{
      \begin{aligned}
        &AB \text{ si } A, B > 0, \\
        &0 \text{ si } A = 0 \lor B = 0
      \end{aligned}
      \right.
    \end{equation*}
  \end{frame}


  \begin{frame}
    \textit{Definición}: Definimos la muliplicación en $\mathbb{R}$ de la siguiente forma

    \begin{equation*}
      AB = \left\{
      \begin{aligned}
        &AB \text{ si } A, B > 0, \\
        &0 \text{ si } A = 0 \lor B = 0, \\
        &-(A(-B)) \text{ si } A > 0 \land B < 0
      \end{aligned}
      \right.
    \end{equation*}
  \end{frame}


  \begin{frame}
    \textit{Definición}: Definimos la muliplicación en $\mathbb{R}$ de la siguiente forma

    \begin{equation*}
      AB = \left\{
      \begin{aligned}
        &AB \text{ si } A, B > 0, \\
        &0 \text{ si } A = 0 \lor B = 0, \\
        &-(A(-B)) \text{ si } A > 0 \land B < 0, \\
        &(-A)(-B) \text{ si } A, B < 0
      \end{aligned}
      \right.
    \end{equation*}
  \end{frame}

  \begin{frame}
    \textit{Teorema}: 

    \begin{itemize}
      \item $-A = (-1)A$. 
    \end{itemize}
  \end{frame}

  \begin{frame}
    \textit{Teorema}: 

    \begin{itemize}
      \item $-A = (-1)A$. 
      \item $AB = (-A)(-B)$.
    \end{itemize}
  \end{frame}

  \begin{frame}
    \textit{Teorema}: 

    \begin{itemize}
      \item $-A = (-1)A$. 
      \item $AB = (-A)(-B)$.
      \item $-(AB) = (-A)B$.
    \end{itemize}
  \end{frame}

\end{document}
